\documentclass[]{article}

\usepackage{amsmath}
\usepackage{multirow}
\usepackage{natbib}
\usepackage{graphicx} 


\begin{document}

% Auto-generated figures file
% \usepackage{graphicx}

\begin{figure}[h!]
    \centering
    \includegraphics[width=0.5\textwidth]{../input_files/plots/step_2_sensitivity_evolution_1739965000.png}
    \caption{Temporal evolution of the partial sensitivities of total energy with respect to mass ($S\_m$, blue) and the damping coefficient ($S\_b$, red) for 20 distinct underdamped oscillators. The plot illustrates a rapid regime shift where the initial dominance of mass-related sensitivity ($S\_m$) transitions to a dominance by damping-related sensitivity ($S\_b$). This crossover, defining the Information Horizon ($T\_H$), consistently occurs early in the transient phase (t < 1s), indicating that the energy model's uncertainty is primarily governed by the damping parameter for the majority of the system's evolution.}
    \label{fig:image1}
\end{figure}

\begin{figure}[h!]
    \centering
    \includegraphics[width=0.5\textwidth]{../input_files/plots/step_2_sensitivity_heatmap_1739965000.png}
    \caption{Temporal evolution of the sensitivity index, S(t), for 20 distinct underdamped oscillators over a 20-second interval. The color intensity corresponds to the sensitivity of the total energy model to perturbations in mass and damping parameters. The plot reveals that sensitivity is highest during the initial transient phase (t < 2s) and subsequently decays, demonstrating that the energy manifold is most susceptible to parameter estimation errors shortly after the initial state.}
    \label{fig:image0}
\end{figure}



\end{document}
                